% Source: https://artofproblemsolving.com/community/c3h2045816_polynomial_and_complex_numbers
% Stable identity: aops-topic-2045816-q01
For real \(r\), the cubic \(P(x)-r\) has real coefficients. It has a non-real root exactly when it does \emph{not} have three real roots (counted with multiplicity), since non-real roots come in conjugate pairs.

\(P'(x)=3x^2+2x-1=(3x-1)(x+1)\), so \(P\) increases on \((-\infty,-1]\), decreases on \(\left[-1,\dfrac{1}{3}\right]\) and increases on \(\left[\dfrac{1}{3},\infty\right)\), with
\[
P(-1)=3,\qquad P\!\left(\dfrac{1}{3}\right)=\dfrac{1}{27}+\dfrac{1}{9}-\dfrac{1}{3}+2=\dfrac{49}{27}.
\]

\begin{asy}
usepackage("amsmath");
size(5.5cm);
defaultpen(fontsize(10pt));
import graph;
real P(real x){return x^3+x^2-x+2;}
draw(graph(P,-2.25,1.25),blue+linewidth(0.9));
draw((-2.5,0)--(1.6,0),Arrow(4)); draw((0,-0.6)--(0,4.3),Arrow(4));
draw((-2.5,3)--(1.6,3),dashed+red); draw((-2.5,49/27)--(1.6,49/27),dashed+red);
label("$y=3$",(1.6,3),E,red); label("$y=\dfrac{49}{27}$",(1.6,49/27),E,red);
dot((-1,3)); dot((1/3,49/27));
label("$x$",(1.6,0),E); label("$y$",(0,4.3),N);
label("\mbox{}",(min(currentpicture,user=true)+max(currentpicture,user=true))/2,black); // keeps the text after the figure black when LaTeX typesets the labels inline
\end{asy}

\begin{itemize}
\item If \(\dfrac{49}{27}<r<3\): by the intermediate value theorem \(P(x)=r\) has a root in each of \((-\infty,-1)\), \(\left(-1,\dfrac{1}{3}\right)\), \(\left(\dfrac{1}{3},\infty\right)\), so all three roots are real.
\item If \(r=3\): \(P(x)-3=x^3+x^2-x-1=(x+1)^2(x-1)\), all roots real.
\item If \(r=\dfrac{49}{27}\): \(P(x)-\dfrac{49}{27}=\left(x-\dfrac{1}{3}\right)^2\left(x+\dfrac{5}{3}\right)\), all roots real.
\item If \(r>3\): on \(\left(-\infty,\dfrac{1}{3}\right]\) we have \(P(x)\le3<r\), and on \(\left[\dfrac{1}{3},\infty\right)\) \(P\) is strictly increasing, so there is exactly one real root, and it is simple because \(P'\ne0\) there. The other two roots are non-real.
\item If \(r<\dfrac{49}{27}\): similarly \(P(x)\ge\dfrac{49}{27}>r\) on \([-1,\infty)\) and \(P\) is strictly increasing on \((-\infty,-1]\), so again exactly one simple real root and two non-real roots.
\end{itemize}
Therefore such a non-real \(z\) exists exactly when
\[
\boxed{r<\dfrac{49}{27}\ \text{ or }\ r>3}.
\]
